\section{Discussion and Limitations}
\label{sec:discussion}

The paper's central question was whether frozen SAM contains texture understanding despite not being built as a texture segmenter. The answer is yes, but only in the sense of \emph{recoverable} texture-relevant structure. The paper does not show that frozen SAM already outputs the right texture partition by default. It shows instead that frozen SAM contains more texture-relevant evidence than its raw outputs reveal.

\paragraph{What we learn from frozen features.} Feature-space exploration answers the existence question positively but cautiously. Coarse frozen features are not texture-blind: they already support nontrivial binary texture partitions under a clustering-only readout. That matters because it shows that the frozen backbone already contains texture-relevant organization before any learned decoder or proposal reasoning.

\paragraph{What we learn from generated masks.} Proposal-/mask-space exploration shows that the generated mask bank also contains substantial texture-relevant structure. On matched RWTD and STLD evaluations, a lightweight readout above frozen proposals recovers strong texture partitions without retraining the proposal generator. The generated masks are therefore not merely noisy by-products; they often already externalize useful texture evidence.

\paragraph{What the two routes jointly imply.} Read together, the routes argue against a simple ``frozen SAM is texture-blind'' conclusion. Frozen SAM carries texture-relevant structure internally and often exposes enough of it in its generated masks to support strong downstream partitions. The difference between the routes is not importance but stage of inspection: features reveal latent organization, while generated masks reveal recoverable partition structure.

\paragraph{Where the routes differ.} The routes are conceptually parallel but empirically asymmetric. The feature probe is reported under an archived probe protocol, so it is most informative about existence and organization of texture structure. The proposal route supports the matched benchmark comparisons and therefore says more about end-task recoverability. Across datasets, RWTD shows fragmented mask evidence that benefits from structured commitment, while STLD often already contains a coherent singleton mask. This is a difference in what the generated masks reveal, not a reason to dismiss either route.

\paragraph{What remains unresolved.} The remaining limitation is not mainly proposal absence. It is ambiguity handling above a strong bank. The current proposal route still makes low-margin wrong-partition commitments on RWTD, and the oracle gap shows that preserving multiple plausible hypotheses for longer could close further headroom. The feature route remains tied to its archived probe protocol, the appendix bridge route does not establish a causal bridge-training benefit, and CAID remains a breadth check rather than a central texture benchmark. The next step is therefore not necessarily a different frozen backbone. It is better uncertainty-aware commitment above fragmented generated masks: calibration, abstention, or multi-hypothesis outputs that avoid premature collapse to the wrong partition.
